% !TEX root = task-4-6-7.tex
% Editable task bank. Append new BankTask blocks at the end of the document.
\documentclass[12pt,a4paper]{article}
\usepackage{physics-textbook}
\usepackage{graphicx,needspace}
\usetikzlibrary{decorations.pathmorphing}
\newif\ifShowSolutions
\ifdefined\TutorEdition
  \ShowSolutionstrue
\else
  \ShowSolutionsfalse % Change to \ShowSolutionstrue to display solutions.
\fi
\newcounter{banktask}
\newsavebox{\BankTaskBox}
\newcommand{\BankTaskContent}[3]{%
  \noindent\textbf{\thebanktask.}\quad #2\par
  \ifShowSolutions
    \Needspace{50mm}% Keep the solution heading with its opening text or diagram.
    #3\par
    \ifdefined\TutorEdition
      \ifcsname TutorSource#1\endcsname
        \nopagebreak\smallskip{\small\csname TutorSource#1\endcsname\par}%
      \fi
    \fi
  \fi
}
\newcommand{\BankTask}[3]{%
  \refstepcounter{banktask}\par\medskip
  \sbox{\BankTaskBox}{\begin{minipage}{\linewidth}%
    \BankTaskContent{#1}{#2}{#3}\end{minipage}}%
  \ifdim\dimexpr\ht\BankTaskBox+\dp\BankTaskBox\relax>\dimexpr\textheight-20mm\relax
    \Needspace{35mm}\BankTaskContent{#1}{#2}{#3}%
  \else
    \Needspace{\dimexpr\ht\BankTaskBox+\dp\BankTaskBox\relax}%
    \noindent\usebox{\BankTaskBox}\par
  \fi
  \medskip
}
\begin{document}
\BookSubsection{4.6.7}{Рассеивающая линза: мнимое изображение предмета}


% Source B81EB6: https://ege.fipi.ru/bank/questions.php?search=1&proj=BA1F39653304A5B041B656915DC36B38&qid=B81EB6
\BankTask{B81EB6}{
В тонкой рассеивающей линзе получено уменьшенное в \(3\) раза изображение предмета. Определите модуль фокусного расстояния линзы, если изображение предмета находится на расстоянии \(f=16\,\text{см}\) от линзы. (в $\text{см}$).
}{
\smallskip\noindent\textbf{Решение от GPT}\par
Изображение в рассеивающей линзе мнимое: \(b=-16\,\text{см}\). Увеличение \(k=|b|/a=1/3\), поэтому \(a=48\,\text{см}\). По формуле тонкой линзы
\[\frac1F=\frac1a+\frac1b=\frac1{48}-\frac1{16}=-\frac1{24}\,\text{см}^{-1}.\]
Следовательно, \(F=-24\,\text{см}\). Ответ: \(|F|=24\,\text{см}\).
}

% Source FDFAF7: https://ege.fipi.ru/bank/questions.php?search=1&proj=BA1F39653304A5B041B656915DC36B38&qid=FDFAF7
\BankTask{FDFAF7}{
Спираль лампочки расположена вблизи главной оптической оси тонкой рассеивающей линзы с фокусным расстоянием, равным по модулю $F$, перпендикулярно этой оси. Расстояние $a$ от линзы до спирали меньше $F$. Затем рассеивающую линзу заменили на собирающую с фокусным расстоянием $F$. Установите соответствие между видом линзы, использовавшейся в опыте, и свойствами даваемого ею изображения.

\TaskMatchingColumns[0.36\linewidth]{Вид линзы}{
А) линза собирающая\par
Б) линза рассеивающая
}{Свойства изображения}{
1) действительное, перевёрнутое, увеличенное\par
2) мнимое, прямое, уменьшенное\par
3) мнимое, прямое, увеличенное\par
4) действительное, перевёрнутое, уменьшенное
}

\begin{center}
\begin{tabular}{|c|c|}
\hline
А & Б \\
\hline
\phantom{3} & \phantom{2} \\
\hline
\end{tabular}
\end{center}
}{
\smallskip\noindent\textbf{Решение от GPT}\par
По формуле тонкой линзы $1/f=1/a+1/b$ для собирающей линзы при $a<F$ получаем $b=aF/(a-F)<0$ и увеличение $\Gamma=-b/a=F/(F-a)>1$: изображение мнимое, прямое, увеличенное (А — 3). Для рассеивающей линзы $f=-F$, поэтому $b=-aF/(a+F)<0$ и $\Gamma=F/(a+F)<1$: изображение мнимое, прямое, уменьшенное (Б — 2). Свойства 1 и 4 не подходят, поскольку оба изображения мнимые и прямые. Ответ: $32$.
}

% Source EFF8E5: https://ege.fipi.ru/bank/questions.php?search=1&proj=BA1F39653304A5B041B656915DC36B38&qid=EFF8E5
\BankTask{EFF8E5}{
Спираль лампочки расположена вблизи главной оптической оси тонкой рассеивающей линзы на расстоянии \(a\) от неё перпендикулярно этой оси, причём \(F<a<2F\), где \(F\) — модуль фокусного расстояния линзы. Затем рассеивающую линзу заменили на собирающую с фокусным расстоянием \(F\). Установите соответствие между видом линзы, использованной в опыте, и свойствами даваемого ею изображения.

\begin{center}
\TaskMatchingColumns[0.36\linewidth]{Вид линзы}{
А) линза рассеивающая\par
Б) линза собирающая
}{Свойства изображения}{
1) действительное, перевёрнутое, увеличенное\par
2) мнимое, прямое, уменьшенное\par
3) действительное, прямое, увеличенное\par
4) мнимое, перевёрнутое, увеличенное
}
\end{center}

\begin{center}
\begin{tabular}{|c|c|}
\hline
А & Б \\
\hline
\phantom{2} & \phantom{1} \\
\hline
\end{tabular}
\end{center}
}{
\smallskip\noindent\textbf{Решение от GPT}\par
Для рассеивающей линзы \(-1/F=1/a+1/b\), поэтому \(b=-aF/(a+F)<0\), а увеличение \(\Gamma=-b/a=F/(a+F)<1\): изображение мнимое, прямое, уменьшенное (2). Для собирающей линзы \(b=aF/(a-F)>0\) и \(|\Gamma|=F/(a-F)>1\), поскольку \(F<a<2F\): изображение действительное, перевёрнутое, увеличенное (1). Вариант 3 невозможен, поскольку действительное изображение предмета в тонкой линзе перевёрнутое; вариант 4 также невозможен, поскольку мнимое изображение здесь прямое. Ответ: \(\boxed{21}\).
}

\end{document}
